% sec-03: outline item C.  Figures 7 (mermaid), 8 (d2), 9 (graphviz).
% The arithmetic spine of the paper.  Tables 8, 9 and 10 must reconcile exactly.
\section{The \$1.6M Gate and the \$3.5M Programme}

\subsection{One Budget, Reused Without Re-derivation}

\draftinstr{Open by stating that the five-year frame is taken verbatim and is
not re-derived here: \$700,000 per year of direct cost, \$3,500,000 over five
years, no cost share, no programme income anticipated, split across four layers
of \$1,600,000, \$720,000, \$780,000 and \$400,000. Source:
\appfile{funding/pdac-funding-applications/final-apply/sections/sec-08-budget-and-leverage.tex}.
Cite \texttt{fundapp2}.}

\subsection{The Award, and What Sits Inside It}

\draftinstr{Explain the difference between an award and the work it buys. An
NIH SBIR award is a total-cost figure: direct costs, indirect costs and fee. Set
Table 8 with the three-component decomposition for both phases, ending at
\$306,000 and \$1,300,000, and state the 15.0 percent overhead-and-fee load
against the 40 percent indirect allowance the company does not claim. Cite
\texttt{nihsbirseed} and \texttt{nihgpsfee}. Source:
\appfile{funding/capitalization-plan/d2/fig-08-two-prices-one-programme.d2.md}.}

\figslot{mermaid-type / state diagram}{6.8}{One award in five states, and the
four guards that must all hold\\at month nine. Two of the four are technical,
one is regulatory,\\and the fourth belongs to an institution the company cannot
bind.}

\draftinstr{After Figure 7, one paragraph per guard, four short paragraphs. G4
gets the longest treatment, because it is the only guard the company cannot
satisfy by working harder, and because the Descoped state exists only for it.
Source: \appfile{funding/capitalization-plan/mermaid/fig-07-phase-gate-state-machine.md}.}

\subsection{Phase I, Nine Months, \$306,000}

\draftinstr{Set Table 9: eight line items summing to \$306,000 total cost. The
rows are personnel for the CEO at 0.6 FTE and the verification engineer at
0.5 FTE, compute, rig hardware, regulatory, site fees, travel, and the indirect
and fee lines. State what Phase I does not buy: no participant is consented and
no drug is dispensed inside it.}

\subsection{Phase II, Twenty-Four Months, \$1,300,000}

\draftinstr{Set Table 10: ten line items summing to \$1,300,000 total cost,
including six participants at \$48,000 each. State the descope path: if the site
executes late, Phase II is resubmitted for three participants rather than six,
which is the Descoped state of Figure 7.}

\subsection{The Same Work Priced Twice}

\figslot{d2-type / layered ledger}{6.6}{One programme, four layers, and two
prices for the same work. The\\SBIR route reaches 1,396,000 of direct effort
inside a 1,606,000\\award; the third column is the 2,104,000 that it does not
reach.}

\draftinstr{After Figure 8, state the reconciliation in prose and give the four
row-level checks. Then explain the delta in one paragraph: 27 months of
additional term, participants 7 to 18, two more dose levels, the correlative
pharmacology, the full VVUQ re-run, and the archive. Do not present the delta as
a shortfall to be apologised for; present it as the honest price of the
mechanism.}

\subsection{What the Money Reaches}

\figslot{graphviz-type / dependency DAG}{7.6}{Seventeen work packages and every
edge between them. Five close\\at 306,000, seven more at 1,606,000, and five are
unreachable\\at that figure, all five downstream of the sixth participant.}

\draftinstr{After Figure 9, one paragraph on the structural finding: the money
runs out one participant after the second dose level clears, and three of the
five unreachable packages are downstream of that point. Then a short paragraph
on WP15 and WP17, the two unreachable packages that are not about participants,
and why WP17's removal changes what \S10 can claim. Source:
\appfile{funding/capitalization-plan/graphviz/fig-09-work-package-dag.gv.md}.}
